\documentclass[UTF8,a4paper,11pt,openany,fontset=fandol]{ctexbook}
\usepackage[a4paper,margin=24mm,headheight=15pt]{geometry}
\usepackage{amsmath,amssymb,bm,mathtools,booktabs,longtable,tabularx,array,graphicx}
\usepackage{enumitem,xcolor,fancyhdr,needspace}
\usepackage[numbers,sort&compress]{natbib}
\usepackage{hyperref,bookmark}
\usepackage{tikz}
\usetikzlibrary{arrows.meta,positioning}
\setCJKmainfont{FandolSong-Regular.otf}[BoldFont=FandolSong-Bold.otf,ItalicFont=FandolKai-Regular.otf]
\setCJKsansfont{FandolHei-Regular.otf}[BoldFont=FandolHei-Bold.otf]
\setCJKmonofont{FandolFang-Regular.otf}
\setmainfont{texgyretermes-regular.otf}[BoldFont=texgyretermes-bold.otf,ItalicFont=texgyretermes-italic.otf,BoldItalicFont=texgyretermes-bolditalic.otf]\setsansfont{texgyreheros-regular.otf}[BoldFont=texgyreheros-bold.otf]
\definecolor{ink}{HTML}{193D55}
\ctexset{chapter={format=\sffamily\huge\bfseries\color{ink},beforeskip=0pt,afterskip=20pt},section={format=\sffamily\Large\bfseries\color{ink}},subsection={format=\sffamily\large\bfseries},subsubsection={format=\sffamily\normalsize\bfseries}}
\hypersetup{colorlinks=true,linkcolor=ink,citecolor=ink,urlcolor=ink,bookmarksnumbered=true,bookmarksdepth=2,pdftitle={Hall MHD与双流体磁重联：模型推导、数值方法与文献调研}}
\setlength{\emergencystretch}{4em}\setlength{\parskip}{3pt}\linespread{1.12}\raggedbottom
\setlist{nosep,leftmargin=2em}\setcounter{tocdepth}{1}\setcounter{secnumdepth}{2}
\pagestyle{fancy}\fancyhf{}\fancyhead[L]{\small\nouppercase{\leftmark}}\fancyfoot[C]{\thepage}
\newcommand{\vct}[1]{\bm{#1}}\newcommand{\dd}{\mathrm d}\newcommand{\VA}{V_{\!A}}\newcommand{\J}{\vct J}\newcommand{\B}{\vct B}\newcommand{\E}{\vct E}\newcommand{\vel}{\vct u}\newcommand{\grad}{\nabla}\newcommand{\R}{\mathcal R}
\newcommand{\note}[1]{\par\smallskip\noindent\textsf{\textbf{学习提示：}}#1\par\smallskip}
\newcommand{\assume}[1]{\par\smallskip\noindent\textsf{\textbf{假设与边界：}}#1\par\smallskip}
\newcolumntype{P}[1]{>{\raggedright\arraybackslash}p{#1}}
\numberwithin{equation}{chapter}

\let\cleardoublepage\clearpage

\tracinglostchars=3

\hypersetup{pdftitle={Hall MHD与双流体模型能否预测离子和电子加热及其分配}}
\newcommand{\loc}[2]{\par\noindent\textsf{原文定位：}\texttt{#1}，#2。\par}
\newcommand{\eV}{\,\mathrm{eV}}
\begin{document}
\frontmatter
\begin{titlepage}
\vspace*{20mm}
{\sffamily\Large Hall MHD 与双流体磁重联调研 · 专题一\par}
\vspace{12mm}
{\sffamily\Huge\bfseries\color{ink}离子与电子怎样被加热\\[8pt]流体模型能预测到哪一步\par}
\vspace{12mm}
{\Large 局地温升、能量分配、不可逆热化与实验约束\par}
\vfill
定向检索与修订：2026年11月13日；11月12日历史批次另列。\\[6pt]
从问题卷原第1.4.5节出发，追踪三个种子的参考文献及后继引用，结合定向检索与原文核验。\\[6pt]
正文区分直接流体计算、动理学对照和实验；不把它们合计为Hall模拟次数。\\[6pt]
这是可追溯的专题证据集，不宣称穷尽全球文献。未取得的版本、勘误、补充及原始数据逐项保留缺口。\\[6pt]
本轮新增全文核读17篇，新增核读主版本共300个PDF物理页；\\[5pt]
三分支完整重读4篇，另由主代理重读F635；新增15个原文PDF（含版本）。\\[5pt]
全仓库全文核读由122增至139篇。

\end{titlepage}
\chapter*{先读结论}
\textbf{可以研究，但必须先选对问题和闭合。}双温Hall MHD、独立物种能量的双流体和十矩模型，已经实际计算离子与电子温度分布、流能和内能预算，并与MAST、HIT-SI及MMS的部分观测比较。它们并非只能研究磁场、完全不能研究加热。直接例证包括Browning等的双温Hall计算、Kaptanoglu等的HIT-SI双温Hall计算，以及TenBarge等的五矩/十矩事件模拟。\citep{L223,Kaptanoglu2020,TenBarge2019MMS}

\textbf{但是，加上Hall项本身不会给出物种热化闭合，更不会自动决定加热比。}Hall项对总电流做功为零，只约束总电磁能量交换的一个分量；各物种仍可受电场做功。磁能进入有序流、成为中心二阶矩内能、相混形成细速度结构后经有限碰撞或明确耗散热化，是相连却不同的步骤；速度空间细化本身不等于精细分布熵产，粗粒化热化须另作定义。双温方程把黏性热给离子、电阻热给电子时，预测依赖这套输运/耗散闭合是否适合所研究的等离子体。

\textbf{最应避免的误读是把所有“比例”看成同一个量。}MRX对称实验的物种电场沉积率约为离子:电子$=2:1$，同一图中的热能变化却约为$1:1$；后者还含储存变化和边界焓流，仍不是局地不可逆熵产。非对称MRX中两物种沉积更接近。温度比、沉积率比、内能变化比和不可逆热化比必须分别报告。\citep{F674,Yoo2017Heating}

\textbf{目前的可信结论是有条件的机制和分布预测。}流体模型可给出下游离子热点、局地电子热点、碰撞交换和热传导的不同时间尺度；Landau/动理学扩展可建立电子相混合至碰撞熵产的链条；十矩可改善压力支撑与部分温度各向异性。但有限矩中的多束、非热尾和非局地热流，跨算例闭合的可靠性，以及同时匹配局地温度、熵产和两物种分配，仍有实质缺口。不存在由本专题证据支持的、与导场、$\beta$、边界和碰撞性无关的普适加热常数。

\textbf{阅读路线：}第1章定位截图和检索范围；第2章区分量并推导能量账本；第3章选择模型；第4章串联流体与动理学论文，并以第\ref{sec:heating_geometry}节逐步解释物理图像；第5章核对实验，第\ref{sec:measurement_chain}节说明诊断和逐图阅读；第6章把已解决与未解决逐项对应到证据，第\ref{sec:gap_map}节区分报告缺口和研究未决；第7章给出可执行的研究设计。末尾的证据索引列出本轮所读版本、页数及详细笔记位置。

本专题是文献研究和选定代数核验，没有运行或复现文中大规模等离子体模拟。前序论文、补充材料、勘误和原始数据只有在实际取得并阅读时才计入范围。所有原文的版权与许可仍按来源页面和原文处理。
\begingroup\setcounter{tocdepth}{2}\small\linespread{1.0}\selectfont\tableofcontents\endgroup
\mainmatter

\chapter{从截图到可追溯的检索范围}
\section{截图原文在哪里}
截图来自修订前\emph{《已解决与未解决的问题》}第1.4.5节，标题“离子加速怎样变成不可逆热化，能量怎样分给两物种”。基线PDF为\texttt{分卷/03\_solved\_open.pdf}的物理第7页、印刷第4页；修订后页码以当前目录为准。维护源是\path{sections/08_solved_open.tex}，通过模型总览章进入合订本，再由分卷脚本生成问题卷。不能只改分卷生成文件。

\begin{longtable}{P{25mm}P{65mm}P{58mm}}
\toprule 稳定键 & 原文与身份 & 原文位置及作用\\\midrule\endfirsthead
\toprule 稳定键 & 原文与身份 & 原文位置及作用\\\midrule\endhead
Yoo2013 & Yoo等，2013，\emph{Observation of Ion Acceleration and Heating during Collisionless Magnetic Reconnection in a Laboratory Plasma}；MRX实验 & \path{papers/Yoo2013.pdf}，5页；图2--4给出电势、流速和下游离子温升。\\
L272 & Yoon与Bellan，2021，\emph{How Hall electric fields intrinsically chaotize and heat ions during collisionless magnetic reconnection}；理论、试粒子与既有PIC解释 & \path{papers/L272_published.pdf}，含封面10页；图1--7；不能当成独立双温Hall比例预测。\\
E512 & Zocco与Schekochihin，2011，\emph{Reduced fluid-kinetic equations for low-frequency dynamics, magnetic reconnection, and electron heating in low-beta plasmas}；KREHM & \path{papers/E512.pdf}，28页；式(76)--(84)、附录A/B。保留电子动理学，不能归为裸Hall。\\
\bottomrule\end{longtable}
表中原文完整书目和DOI见参考文献。\citep{Yoo2013,L272,E512}

\section{怎样追引，哪些地方没有覆盖}
检索日期为2026年11月12日。以三个种子的DOI定位OpenAlex记录，分别取得参考文献及引用它们的记录。Yoo2013得到34条后向引用和69条前向引用，L272得到49和6条，E512得到95和126条；共379条有向关系、363个去重节点。三条种子的一跳查询完成；数据库的引用图不等于全球引用全集，也没有递归穷尽每一代后继。

宽检索式包括\texttt{Hall MHD ion electron heating}、\texttt{two fluid reconnection heating}、\texttt{two temperature Hall MHD}，以及EMHD、ten moment、Landau fluid、gyrofluid和energy partition等组合。前两式分别取得5133/5133、9464/9464条，第三式取得3800/15993条后遇到HTTP 429；其余七式首轮也受限。共保存16020个去重元数据记录，\textbf{其中有大量Hall传热、工程纳米流体等不相关命中，不能称为16020篇相关论文，也未逐条人工筛完。}限流不是零命中。

人工优先追踪三个种子的原文参考文献、直接相关后继，并补充双温Hall、十矩、MAST/HIT-SI/MRX及热流闭合的定向网页检索。原文优先取出版机构、作者/机构仓库、arXiv和会议官方网站。详细检索式、查询状态、来源、原始边和下载日志保存在本目录\path{search/}；可重跑脚本是根目录\path{scripts/search_ion_heating.py}。题名分流不等于全文相关性核验。

取得可解析PDF后，核对题名、作者、版本；各分支阅读全文、方法、结果、讨论、文内附录及参考文献，直接检查关键公式和图表。详细记录见\path{reading_notes/}的三个分支笔记及JSON。旧论文重新阅读不增加独立全文数量，预稿与正式稿也不重复计数。

\section{仍应追踪的具体缺口}
11月13日新增核读Drake等2009年JGR反平行拾取论文（Drake2009Pickup，DOI 10.1029/2008JA013701），14页含附录A。历史下载失败实际指向另一篇2009年ApJ强导场/重离子丰度论文（DOI 10.1088/0004-637X/700/1/L16）；后者仍未核读，不能用前者填补其阈值证据。2003年EMHD系列第三篇电子加热、Stanier学位论文等也仍是待核读连接点，不据摘要作科学结论。Kaptanoglu2020存在DOI为\texttt{10.1063/5.0032632}的勘误，本轮未取得勘误正文；因此本专题不以未经勘误对照的精确温度比例作核心结论。TenBarge2019的正式网页比已读arXiv版多一个图及数值方法节，正式PDF访问受限；已读版本及网页补查分别记录。

这不是以“搜集所有”为由虚构穷尽性：当前可交付的是明确检索路径、取得的原文、完成的全文证据和可继续执行的缺口清单。未读候选可在\path{candidate_screening.csv}及检索日志中反查，不能作为已解决问题的证据。

\chapter{先分清加速、温升、热化和比例}
\section{温度是什么：一个两束反例}
设离子分布为$f_i(\vct x,\vct v,t)$，电子分布为$f_e$。对单电荷离子，令准中性数密度$n_i=n_e=n$，$e>0$，质量分别为$m_i,m_e$。定义
\begin{align}
\vct u_i&=n^{-1}\int\vct v f_i\dd^3v,&\vct u_e&=n^{-1}\int\vct v f_e\dd^3v,\\
\vct P_i&=m_i\int(\vct v-\vct u_i)(\vct v-\vct u_i)f_i\dd^3v,&
\vct P_e&=m_e\int(\vct v-\vct u_e)(\vct v-\vct u_e)f_e\dd^3v.
\end{align}
两个向量并置表示外积，$P_{i,ab}$是离子压力张量的$a,b$分量。标量压力$p_i=\operatorname{tr}\vct P_i/3$、$p_e=\operatorname{tr}\vct P_e/3$。本文温度以能量计，$T_i=p_i/n,T_e=p_e/n$；换成开尔文要除以$k_B$，换成eV要除以基本电荷的数值。内能密度与有序动能密度为
\begin{equation}
 U_i=\tfrac12\operatorname{tr}\vct P_i=\tfrac32nT_i,\quad
 U_e=\tfrac32nT_e,\qquad K_i=\tfrac12m_inu_i^2,\quad K_e=\tfrac12m_enu_e^2.
\end{equation}
这一定义对非Maxwell分布也成立。温度张量或中心二阶矩变大，不自动证明分布已经热化。此处共用$n$是准中性写法；讨论全Maxwell双流体时，应恢复独立$n_i,n_e$及电荷密度，不能把准中性当作完整Maxwell的恒等约束。

\textbf{自设教学反例，非论文模拟：}两束等密度离子分别以$+V\hat x,-V\hat x$漂移，每束各向同性温度均为$T_0$。混合后平均流速为零，
\begin{equation}
 T_{xx}=T_0+m_iV^2,\quad T_{yy}=T_{zz}=T_0,\quad T_i=T_0+\frac{m_iV^2}{3}.
\end{equation}
仪器若只辨认混合线宽，会得到较高二阶矩温度；但两束各自并没有变热。碰撞可以把这份相对漂移能重新分配而形成近Maxwell分布。这正是Yoon2023相空间分析和Yoo2014无/有碰撞对照要求注意的区别。\citep{Yoon2023PhaseSpace,F679}

\section{分别写出两物种的能量账本}
以下是无电离、辐射、外加粒子源时的非相对论矩方程；实际装置如有这些源必须增加，不能让数值残差代替它们。先保留完整压力与热流，不预先假定各向同性：
\begin{align}
 \partial_t U_i+\nabla\!\cdot(U_i\vct u_i+\vct q_i)
 &=-\vct P_i:\nabla\vct u_i+H_i,\label{eq:ioninternal}\\
 \partial_t U_e+\nabla\!\cdot(U_e\vct u_e+\vct q_e)
 &=-\vct P_e:\nabla\vct u_e+H_e.\label{eq:electroninternal}
\end{align}
这里$\vct q_i,\vct q_e$是相对各自平均速度定义的热流，单位$\mathrm{W\,m^{-2}}$；$H_i,H_e$是碰撞及所选物理热源的内能贡献，单位$\mathrm{W\,m^{-3}}$。双点积按$\vct P:\nabla\vct u=\sum_{a,b}P_{ab}\partial_bu_a$定义。方程尚未闭合，压力张量和热流需要各自演化或本构关系。

离子和电子连续性分别为$\partial_t n+\nabla\cdot(n\vct u_i)=0$与$\partial_t n+\nabla\cdot(n\vct u_e)=0$。把各自动量方程点乘速度，再用连续性消去密度导数，可得
\begin{align}
\partial_tK_i+\nabla\!\cdot(K_i\vct u_i+\vct P_i\cdot\vct u_i)
 &=en\vct u_i\cdot\E+\vct P_i:\nabla\vct u_i+\vct u_i\cdot\vct R_i,\\
\partial_tK_e+\nabla\!\cdot(K_e\vct u_e+\vct P_e\cdot\vct u_e)
 &=-en\vct u_e\cdot\E+\vct P_e:\nabla\vct u_e+\vct u_e\cdot\vct R_e.
\end{align}
用了恒等式$-\vct u\cdot\nabla\cdot\vct P=-\nabla\cdot(\vct P\cdot\vct u)+\vct P:\nabla\vct u$，以及$\vct u\cdot(\vct u\times\B)=0$。$\vct R_i,\vct R_e$为单位体积碰撞动量交换。于是压力功在流能和内能之间成对抵消；电场功首先进入总粒子能量账本，并不全是即时热源。

再加电磁Poynting定理
\begin{equation}
\partial_t\left(\frac{\epsilon_0E^2}{2}+\frac{B^2}{2\mu_0}\right)
+\nabla\cdot\frac{\E\times\B}{\mu_0}=-\J\cdot\E,\qquad
\J=en(\vct u_i-\vct u_e),
\end{equation}
离子和电子电功之和恰为$\J\cdot\E$。全Maxwell模型保留电场能，低频Hall约化按其排序略去相应小量，但磁能和粒子能量预算仍须一致。守恒检查要求
\begin{equation}
\vct R_i+\vct R_e=0,\qquad H_i+H_e+\vct u_i\cdot\vct R_i+\vct u_e\cdot\vct R_e=0.
\end{equation}
上述零和条件仅适用于封闭物种碰撞交换；若$H_i,H_e$包含外加热、辐射等，右端须保留相应净输入/输出。\textbf{不能只给总能方程加电阻热，却在另外两套内能方程重复加一次。}十矩模型的内能已经是压力迹，也不能再独立推进一个未经协调的标量压力。

\section{Hall总功为零为什么仍容许物种增能}
在离子速度形式的广义欧姆定律中，Hall电场分量写作
\begin{equation}
\E_H=\frac{\J\times\B}{en},\qquad \J\cdot\E_H=0.
\end{equation}
但是分别点乘物种电流，有
\begin{equation}
en\vct u_i\cdot\E_H=\vct u_i\cdot(\J\times\B),\qquad
-en\vct u_e\cdot\E_H=-\vct u_e\cdot(\J\times\B).
\end{equation}
二者可非零，和才为零，因为$\vct u_i-\vct u_e=\J/(en)$。这是特定欧姆分解中的恒等关系，\textbf{不能把实验全部面内电场一律命名为纯Hall项}：压力梯度、惯性和其他分量可能同时贡献。Hall项也不自动破坏电子磁冻结；实际破冻结仍须检查电子压力张量、惯性、电阻或数值项。

因此“Hall项不直接消耗总磁能”与“Hall重排场、流和电流并改变物种能量路径”相容。L263所谓Hall电流电阻耗散，是改变后的电流经$\eta J^2$损失；其压力绝热演化未回收这份热，更不能拿它的比例作为物种升温预测。\citep{L263}

\section{电阻、黏性和热传导各自做什么}
取简化各向同性摩擦$\vct R_e=\alpha(\vct u_i-\vct u_e)$、$\vct R_i=-\vct R_e$，$\alpha\geq0$。则总流能碰撞损失为
\begin{equation}
\vct u_i\cdot\vct R_i+\vct u_e\cdot\vct R_e
=-\alpha|\vct u_i-\vct u_e|^2=-\eta J^2,\quad
\eta=\frac{\alpha}{e^2n^2}.
\end{equation}
这里$\eta$为SI电阻率，单位$\Omega\,\mathrm m$，不是磁扩散率$\eta/\mu_0$。总内能因此应增加$\eta J^2$；守恒只限定两物种之和，\textbf{具体给谁多少仍需碰撞算子/质量比排序}。若用$C(T_e-T_i)$表示电子向离子的温差交换，离子取正、电子取负；它只交换总内能，不能同时向两物种增热。某些简化模型把主要电阻热给电子、主要黏性热给离子，这是可说明且可检验的近似，不是Hall恒等式的推论。

将压力写为$\vct P_i=p_i\vct I+\vct\Pi_i$，则
\begin{equation}
-\vct P_i:\nabla\vct u_i=-p_i\nabla\cdot\vct u_i-\vct\Pi_i:\nabla\vct u_i,
\end{equation}
电子同理。第一项是压缩/膨胀；绝热压缩也能提高温度。若牛顿型黏性关系为$\vct\Pi_i=-\vct\tau_i$且输运系数满足正定条件，第二项才可与正定黏性耗散对应；对一般无碰撞压力张量，它可以双向交换，不能逐像素称作不可逆熵产。

热流散度能使局地温度升降，却主要搬运能量。对于标量近Maxwell模型、物理温度$\theta_s=T_s/k_B$、单位体积熵$\mathcal S_s$，其熵预算含
\begin{equation}
\partial_t\mathcal S_s+\nabla\cdot\left(\mathcal S_s\vct u_s+\frac{\vct q_s}{\theta_s}\right)
=\frac{Q_{s,\mathrm{diss}}}{\theta_s}+\vct q_s\cdot\nabla\frac{1}{\theta_s}
+\hbox{物种热交换贡献}.
\end{equation}
Fourier热流$\vct q=-\kappa\nabla\theta$给出正的传导熵产，但局地$-\nabla\cdot\vct q$可正可负。每种粒子的热交换熵项也不必各自为正；两物种总和满足第二定律。强非Maxwell时应使用分布熵和对应碰撞模型，不能机械套标量熵公式。

\section{四种比例必须用四个名字}
对固定控制体$V$、时间窗$[t_0,t_1]$定义：物种电场沉积$W_i=\int\!\dd t\int_Ven\vct u_i\cdot\E\,\dd V$、$W_e=\int\!\dd t\int_V-en\vct u_e\cdot\E\,\dd V$；储存内能变化$\Delta\mathcal U_i=\int_V[U_i(t_1)-U_i(t_0)]\dd V$，电子同理。它们与不可逆热化$Q_i^{\rm irr},Q_e^{\rm irr}$以及温度之比是不同量：
\begin{equation}
R_W=\frac{W_i}{W_e},\qquad R_U=\frac{\Delta\mathcal U_i}{\Delta\mathcal U_e},\qquad
R_Q=\frac{Q_i^{\rm irr}}{Q_e^{\rm irr}},\qquad R_T=\frac{T_i}{T_e}.
\end{equation}
$Q_s^{\rm irr}$必须依所选模型另给操作定义，例如Numata2015所定义的碰撞自由能耗散率$D_s$的时空积分，或近Maxwell输运模型中明确定义的耗散热；本专题不假定各种定义天然等价。一般熵产不能统一换算成新增内能，传导和零迹各向同性化就说明二者不同。若分母接近零或量可变号，不能把比值包装成稳定的“份额”。局地、全域、瞬时和时间积分结果也必须分开。由式\eqref{eq:ioninternal}直接积分得
\begin{equation}
\Delta\mathcal U_i=\int_{t_0}^{t_1}\!\dd t\int_V(-\vct P_i:\nabla\vct u_i+H_i)\dd V
-\int_{t_0}^{t_1}\!\dd t\oint_{\partial V}(U_i\vct u_i+\vct q_i)\cdot\dd\vct A.
\end{equation}
电子亦然。密度、驻留时间、热流和边界焓输运改变温升；绝不能只从一张$T_i/T_e$图反推出当地磁能的不可逆分配。严格无碰撞Vlasov演化保持精细分布熵；有限矩、有限分辨率或粗粒化熵增加的含义必须交代。

\chapter{选哪一种模型才能回答哪一种问题}
\begingroup\small
\begin{longtable}{P{31mm}P{57mm}P{60mm}}
\toprule 实际模型 & 能计算什么 & 对物种热化比例的边界\\\midrule\endfirsthead
\toprule 实际模型 & 能计算什么 & 对物种热化比例的边界\\\midrule\endhead
等温Hall/EMHD & 磁场、流速、电势及指定温度下的响应 & 温度被规定，不能宣称独立预测升温。固定离子EMHD没有自洽离子能量。\\
单总压或固定$T_i/T_e$ & 总流体热力学、流场及Hall结构 & 分配是输入假设；Stanier2013取$p_i=p_e=p/2$，不能输出新的加热比。\\
两物种独立标量绝热方程 & 两种温度随压缩、对流改变 & 平滑理想过程可逆；若靠网格耗散升温，需说明并检验。\\
双温耗散Hall MHD & 分别的温度、热流、黏性/电阻热与物种交换 & 可预测所选闭合内分配；需要检验输运、异常电阻/超电阻及损失边界。\\
五矩双流体 & 两套密度、动量、标量能量；可搭配全Maxwell或准中性 & 各向同性、热流闭合限制；解出能量方程不保证无碰撞热化正确。\\
十矩双流体 & 两套完整对称压力张量；各向异性和非旋转对称压力 & 仍需三阶热流闭合；不能唯一恢复多束、非热尾、俘获与速度空间熵。\\
Landau/陀螺流体 & 有限矩近似平行相混、部分FLR和非局地响应 & 依赖排序及闭合；L244实际采用冷离子，不能给两种热化竞争。\\
KREHM & 流体变量加电子漂移动理学/Hermite高阶矩 & 能研究排序内电子熵产；E512省去的离子通道使其不是普适两物种比例模型。\\
Hall/EMHD加试粒子 & 给定流体场下的轨道、加速与速度展宽 & 无自洽反作用及竞争能量预算；测试粒子能量不能任意加回守恒总能。\\
动理学离子加流体电子 & 自洽离子分布及电子闭合反馈 & 属混合模型，电子热化仍受闭合限制；Finelli2021未完成详细能量分配研究。\\
双物种GK/PIC & 更丰富的相空间和粒子/场耦合 & 仍受质量比、碰撞模型、分辨率、统计噪声、排序及边界限制；并非天然的完美真值。\\
\bottomrule\end{longtable}
\endgroup
以上边界分别由所读方法与实际算例确定。\citep{Stanier2013,L223,Kaptanoglu2020,L247,TenBarge2019MMS,L244,E512,Yoon2018Heating,Finelli2021,Numata2015Heating}

\section{为什么完整矩方程不等于有限闭合已经可靠}
动理学矩层级能严格写出，但压力演化需要三阶热流，热流又需要更高阶矩。两种分布可以有相同密度、速度和压力，却有不同热流、尾部和后续响应；有限数量的低阶矩不能唯一辨认它们。这是信息不足，而非守恒恒等式错误。

以常用十矩局地闭合为例：$\partial_kQ_{ijk}=\nu(P_{ij}-p\delta_{ij})$。取迹得$\nu(\operatorname{tr}P-3p)=0$，所以该松弛\textbf{不直接增加物种总内能}，但会改压力各向异性、场和压力功，从而改变后续分配。它也不能替代真正跨空间输运的热流。TenBarge2019和后续梯度闭合的差异应归于闭合与模式响应，而不是一句“双流体能量不守恒”。\citep{TenBarge2019MMS,Ng2020,Closure2026}

多群电子十矩F635进一步说明这一点：在固定离子背景、两电子群条件下分开演化可改善反向流群的表达，但并未演化竞争的离子加热；群内部仍可能存在反向束，密度下限和高阶矩缺失也继续存在。它不能作为离子/电子比已经解决的证据。\citep{F635}

\chapter{论文实际把问题推进到了哪里}
\section{直接双温Hall：从固定温比到局地热点}
\subsection{MAST：Stanier2013到Browning2016}
Stanier2013的岛合并Hall计算解释密度双峰、几何压缩、磁通堆积及反弹，但取$p_i=p_e=p/2$；其关于未来独立温度的讨论不能倒写成已经预测了分物种加热。\citep{Stanier2013}

Browning等L223综述并扩展这条模拟线，图3报告初步双温环形Hall计算：离子因出流黏性耗散在下游形成双峰，电子耗散位置不同。典型归一化取$n_0=5\times10^{18}\,\mathrm{m^{-3}}$、$L_0=1$ m、$B_0=0.5$ T，初始$T_i=T_e=10\eV$，$d_i/L_0=0.145$；无量纲$\eta=10^{-5}$、$\nu_i=10^{-3}$，电子平行导热和离子垂直导热系数分别$0.1$和$10^{-7}$，图3超电阻$\eta_H=10^{-8}$。初始正文的总压力式本身没有两物种分配，双温扩展在后文另述。\citep{L223}

这里首次需要正面回答用户的问题：\textbf{该模型确实能研究局地离子与电子温升及不同位置，不能说Hall完全做不到。}但电子热点由自由超电阻参数强烈控制；双温方程未完整印出，图3的网格、时间步及独立收敛序列未报告。离子峰可达keV量级，与较低驱动实验并非逐点定量吻合。另一个Taylor松弛模型估计的是可释放磁能及其近似$B_p^2$标度，作者明确不由它决定物种分配和空间分布。
\loc{L223}{16页作者稿全部；PDF第3--4页模型，第5--7页双温与图3，第9--14页松弛比较}

\subsection{HIT-SI：三维双温Hall及装置损失}
Kaptanoglu等用PSI-Tet研究HIT-SI驱动装置，把离子与电子能量分别演化，比较单温和双温，结合干涉、离子Doppler、Thomson散射及磁探针约束。模型能区分黏性/压缩驱动的离子升温、电子热输运、碰撞耦合及壁面损失。双温选择会反馈黏性、驱动电流和整体平衡，因此局地加热不只是给既有磁解事后贴一张温度图。\citep{Kaptanoglu2020}

其价值是把双温Hall推进到三维受驱动实验几何和多诊断比较；它不是一张孤立反平行电流片的普适重联加热规律。PSI-Tet用103887个三阶四面体单元，约9 mm等效分辨率；增强电子质量使质量比为100，电子惯性长度约12 mm。Crank--Nicolson步长最大40 ns，高频约10 ns；扫描14.5--68.5 kHz驱动。NIMROD还采用不同的注入器边界近似，比较并非只切换温度方程。图3多为含注入器全体积预算，图4两种温度大尺度空间结构相似、幅度不同，不能称两者热点全然异位。

缺少真实电路模型使作者明确不作严格定量实验验证；中性粒子、回收、辐射及更真实热边界仍是缺失物理，粒子扩散也未达到参数不敏感极限。已读35页接收稿与较早arXiv稿题名不同，按DOI识别为同一工作；存在未取得正文的勘误，本文因此只采用结构性发现和明确证据边界，不以精确比例作为已验证定律。
\loc{Kaptanoglu2020}{35页接受稿全部；方程、温度/功率图及实验比较详见分支笔记，勘误缺口见检索日志}

\section{双温输运与激波：可以追踪热化，但不等于Hall层已经解决}
Longcope与Bradshaw2010用一维准中性、共同流速但独立温度的碰撞流体描述倾斜重联后慢激波与传导前沿。离子先经黏性子激波加热，电子随后通过碰撞交换和传导响应；电子热前沿与离子温度峰的位置及时间并不重合，常见单温假设会改变热结构。该例展示两温方程对加热因果顺序的实际作用。\citep{L164}

然而共同流速使它不是自洽求Hall电场、电子扩散区和重联率的模型。沿场近一维几何、经典碰撞输运、热流适用性及离子子激波尺度均有限制；若平均自由程与梯度尺度不分离，需要非局地/动理学处理。原文没有一个跨所有重联环境的$Q_i/Q_e$。图像不能与弱碰撞MRX原封不动互换。
\loc{L164}{45页全文；第10--12页式(4)--(7)、(14)--(18)，图4--7、9--14}

Li与Livescu2019的L247则给出含完整输运的高阶双流体数值框架：采用准中性、忽略电子惯性和位移电流，一般方程保留独立物种内能、黏性、热流和碰撞交换，足以作为实现起点；它不是未约化的全Maxwell双流体。\textbf{但第IV节的四个验证算例强制$T_i=T_e$并删去离子能量方程。}因此其数值精度、波/流验证不能被宣传成离子/电子加热比例已验证。方程能力与实际算例能力必须分栏登记。\citep{L247}
\loc{L247}{44页全文；PDF第16页的验证简化条件，以及后续四个验证问题}

\section{Hall电场到离子轨道：2018、2019、2021、2023的递进}
Yoon与Bellan2018在恒密度固定离子背景EMHD中产生场，再另积分测试离子。二维Harris片约$200\times600$网格、$30\times120\,d_e$，6000个初始静止测试离子，示例质量比10并比较更大质量比的定性行为。增长电场使近邻轨道分离、速度分布展宽，给出电场梯度失磁化路径。它比“有Hall场所以热”更具体，但测试离子不反馈流体场，未演化电子与离子竞争预算，也没有碰撞Maxwell化。\citep{Yoon2018Heating}

2019年F743改用自洽PIC：双Harris片、$m_i/m_e=100$，$1024\times2048$网格覆盖$40.96\times81.92\,d_i$，导场比$0,0.1,0.3,0.5$。电场梯度指标与高离子温度区相关，较强导场抑制这套机制，重离子响应也有差异。静态Harris平衡导出的导场阈值不能直接套到实时重联；部分平行温升区并未由同一指标解释。\citep{F743}

2021年L272建立更统一的场梯度/正则涡量说明，图6--7延续2019数据，而非又一次独立验证。给定场试粒子的$T_E$是\emph{能量分布标准差}，不等于标准中心速度温度。反抛物势中的指数轨道分离也不足以单独证明有界混沌或熵增；原文系数疑点在原库讲义中已核验。\citep{L272}

2023年Yoon、Bellan与Yun进一步用相空间分解说明：总$T_{xx}$可以因两束反向流增大，而单束温度低于上游；有序束混合与混沌展宽应分开。二维反平行PIC采用$1024\times8192$、$51.2\times409.6\,d_i$、质量比100、每格200--1000粒子。冻结PIC时刻场上用Benettin方法计算有限时间Lyapunov指数，比单看两轨道更直接，但仍不等于全时变三维证明。\textbf{作者明确留下磁能进入离子的精确份额，以及有序与随机增能相对份额。}\citep{Yoon2023PhaseSpace}
\loc{Yoon2018Heating / F743 / L272 / Yoon2023PhaseSpace}{依次7、5、10、9个PDF物理页全部；重点为2018图1--3，2019图1--5，2021图1--7，2023图1--7及附录}

因此这条链解决的是一条更可检验的离子增能机制及其条件，不是“纯Hall热方程已能无参数输出所有物种热化比”。MRX与这些图像之间应保留原始观测和后续机制解释的时间顺序。

\section{电子相混：从KREHM到Landau有限矩，再到双物种参照}
E512的低频、强导场、低$\beta_e\sim m_e/m_i$排序保留电子平行相混合、电子惯性与特定离子FLR响应；高阶电子分布$g_e$一般不能始终置零。流体自由能经可正可负的耦合进入高阶矩，最终由正定碰撞熵产耗散；自由能增加本身不是热量。离子Landau阻尼和熵级联被排序掉，因此它没有竞争性的完整离子热化。\citep{E512}

Loureiro等2013非线性KREHM计算直接验证电子耗散的时空滞后：主例重联率峰约$22.7\tau_A$，电子耗散峰约$29.4\tau_A$，热化集中于分离线附近；电子高阶Hermite谱连接到耗散。模型使用Hermite截断及超碰撞/空间超扩散；文中“大系统约60\%初始磁能转为电子热”既未在相应图中展示，也不是两物种电子份额60\%。\citep{Loureiro2013Heating}

Tassi等L244把Landau响应压缩进有限矩陀螺流体热流闭合；磁能经电子温度扰动能耗散，斜撕裂可变慢或受抑。但实际非线性算例采用冷离子，离子温度扰动为零。其耗散峰没有上述动理学滞后，作者主动质疑是否过度耗散、是否需更多矩；无耗散运行的数值损失虽仅为初始二次自由能$H(0)$的约0.002，却已约为释放磁能的五分之一，不能用背景总热能作分母掩盖。\citep{L244}

Numata与Loureiro2015提供重要的\emph{双物种陀螺动理学参照}：用碰撞算子定义$D_i,D_e$的不可逆耗散率，在$\beta_i=\beta_e=0.01$至1的扫描中，低$\beta$电子主导，高$\beta$离子相混作用变强，离子偏磁岛内部、电子偏分离线/边缘。$\beta=1$时部分晚期瞬时$D_i/D_e$达量级1，并非累计热化普适相等。扫描同时改变$d_e/a$，碰撞率人为同设正则化，质量比100，不能包装成真实参数的单变量实验。\citep{Numata2015Heating}
\loc{E512 / Loureiro2013Heating / L244 / Numata2015Heating}{28、6、27、18页全部；E512式(76)--(84)，Loureiro图1--5，L244图4及结论，Numata式(2.5)、图5--7、11--12与附录A}

这些论文把“热化”从温度图推进到自由能和熵预算，但有的离子通道被删除，有的已不是有限矩流体。不能把它们的进展全部归功于Hall项，也不能因为需要动理学便否认有限矩在适用范围内的价值。

\section{十矩、混合闭合与局地湍流：哪里有效，哪里失败}
TenBarge等2019用全Maxwell五矩/十矩模拟2015年10月16日MMS事件。十矩二维和三维的全域流能/热能预算相近，图1中离子热能增长大于电子；这直接证明有限矩可以研究分物种能量。但五矩中LHDI过强并导致显著能量丢到网格，不能用这组运行的比例作物理结论。十矩固定$k_0d=0.1$闭合又把LHDI压得过强。电子$T_\parallel/T_\perp$峰约2.5接近观测量级，但电子各向异性、平行优先升温区在磁鞘侧，MMS与对应PIC位于磁层侧；该侧别问题不等于所有总温度热点都错位。\citep{TenBarge2019MMS}

原库Ng2020与Closure2026的梯度闭合改进了离子层、尺度效应和部分三维不稳定性，但热流、电子各向异性侧别与跨场扩散仍是限制。改进闭合能改善结果，说明这些误差不能全称为“双流体原理不允许”；但保留更多矩也没有自动完成不可逆热化验证。\citep{Ng2020,Closure2026}

Finelli2021比较动理学离子加等温电子、动理学离子加Landau流体电子和PIC。电子闭合改善后，电子喷流、各向异性和离子反馈更接近PIC；俘获、电子FLR和部分不稳定性仍缺失，作者明确把详细能量分配留待未来。这是电子闭合如何反馈离子的证据，不能写成已测得$Q_i/Q_e$。\citep{Finelli2021}

Che2014标题虽含“Two-Fluid”，实际先做三维PIC的Buneman湍流诊断，再建立近似流体描述；原文明确\textbf{没有发生重联}。其电子平行/垂直能量比不是离子/电子比。模型需另给动理学理论或PIC标定的拖曳$D_{ez}=f(t)$，尚非独立可迁移的闭合；约化分布熵的变化也需交代粗粒化。\citep{Che2014Buneman}
\loc{TenBarge2019MMS / Finelli2021 / Che2014Buneman}{17、17、26页作者稿全部；TenBarge图1、6--10，Finelli图3--8，Che图1、6--9及式(7)--(19)}

% 2026-11-13：与既有逐篇结果配套的因果讲解；原文图号保持作者编号。
\section{先在空间里看见能量路径}
\label{sec:heating_geometry}
下面的图是本报告绘制的\textbf{物理示意，不是实验数据或计算结果}。用$N$表示入流法向、$L$表示出流、$M$表示面外电流/导场方向，以免混淆不同论文的坐标：MRX中$(N,L,M)=(R,Z,Y)$；Yoon系列中为$(x,y,z)$；Drake2009中为$(y,x,z)$。MAST两环沿$z$合并，出流主要沿$r$，导场为环向。局部几何相似并不保证边界、碰撞率与加热机制相同。
\begin{figure}[htbp]
\centering
\begin{tikzpicture}[>=Stealth,font=\small]
\fill[blue!5] (-6,-2.3) rectangle (6,2.3);
\draw[ink,thick] (-5.8,1.9) .. controls (-2,1.8) and (-1,0.1) .. (0,0)
 .. controls (1,-0.1) and (2,-1.8) .. (5.8,-1.9);
\draw[ink,thick] (-5.8,-1.9) .. controls (-2,-1.8) and (-1,-0.1) .. (0,0)
 .. controls (1,0.1) and (2,1.8) .. (5.8,1.9);
\draw[->,thick] (0,2.2)--(0,0.7);
\draw[->,thick] (0,-2.2)--(0,-0.7);
\draw[->,thick] (0.7,0)--(3.4,0);
\draw[->,thick] (-0.7,0)--(-3.4,0);
\fill[orange!25] (0,0) ellipse (0.6 and 0.35);
\node at (0,0) {X};
\node[align=center] at (-4.6,0) {下游离子\textbf{温升}\\停滞、混合、耗散};
\node[align=center] at (4.6,0) {下游离子\textbf{温升}\\需要另查分布};
\node at (1.5,2.0) {入流 $N$};
\node at (2.1,-0.45) {出流 $L$};
\node[align=center,fill=white] at (-3.4,1.55) {分离线附近：电势、\\强梯度、物种解耦};
\node[align=center,fill=white] at (3.5,-1.65) {电子热可沿场输运；\\热化可滞后于重联};
\node at (0,-2.7) {$M$垂直纸面；局地电功峰、温度峰、熵产峰无需重合};
\end{tikzpicture}
\caption{将Yoo2013/F679、Tanabe2015MAST和电子相混研究放入共同几何的阅读示意。橙区只标X点邻域，不代表所有装置的电子热点宽度；强导场MAST、零导场MRX和强导场KREHM的具体图像须分别核对。}
\end{figure}

先看重联前后磁张力如何产生出流，再看电场如何给各物种做功，最后检查流能转为压力、热流把能量带往哪里，以及何种过程消去细速度结构。这三层各有证据：场/流图约束第一层，谱线或分布的中心矩约束第二层，碰撞算子、熵预算或分布趋向Maxwell的证据约束第三层。只有温度图，不能跳过中间两步。

\subsection{双温Hall的热点：磁张力释放以后，热在哪里留下来}
Browning等（2016，\texttt{L223}）的两环合并中，环向电流形成极向重联场；重新连接的场线张力推动径向出流。离子出流在下游减速、剪切，指定的离子黏性把有序流能变成内能，因而在X点两侧形成离子双峰。电子则通过电阻和超电阻通道增热，沿场热输运又改变其温度分布。图2先读磁通和电流层怎样挤压、出流、形成小岛，再读图3的温度随时间变化，才能明白两种峰为什么不同。\citep{L223}

这里有一个容易被概括丢掉的细节：\textbf{该文电子温度由超电阻耗散强烈控制，不能简单称为“电流最大处的欧姆热点”。}超电阻与电流梯度有关，作者明确说温度峰与电流密度不必同位；中心有峰或空心可能与小岛形成、抛出有关。这是指定耗散模型的机制解释。MAST实测X点窄峰是另一项证据，二者在空间上具有联系，并不构成对超电阻系数的独立测量。图3仍是初步双温结果，不能从其峰高倒推出已验证的局地$Q_i/Q_e$。
\loc{L223}{PDF第4--7页图2、3及第5--6页双温结果；第3页式(1)--(5)先给总压力模型，不能冒充完整打印的双温扩展}

\subsection{碰撞激波的因果顺序：离子先热，电子把热送到上游}
Longcope与Bradshaw（2010，\texttt{L164}）研究重联后磁通管的沿场对撞/减速。图4右图的流动从右向左，沿途依次是电子传导前沿CF、离子子激波、离子冷却区IC。离子在子激波由超声速降至亚声速，黏性应力做功使$T_i$猛升；电子因质量小、携带的流动能少，同步减速却不立即取得同样热量。下游离子通过较慢的物种碰撞交换把能量给电子，电子高平行热导率又将热传回子激波前方，形成预热区。电子压强梯度对应的极化电场还会提前减速离子，故“前沿”并非一块与流动无关的热斑。\citep{L164}

图4示例$M_1=3.5$、入流马赫数$M_i=2.4$，两者不是重联率或Alfvén马赫数。离子Prandtl数$0.33$，电子$Pr_e^{-1}=77$，解释了窄离子子激波与长电子前沿；图5比较三段空间尺度，图9--11说明它们怎样逐渐形成。因此同一位置先看到$T_i>T_e$、晚时趋同，既可能是能量交换和输运，也不能被误解为Hall项在不同时刻“换了加热对象”。该模型用共同体流、经典碰撞输运，不能自行检验MRX无碰撞层的轨道混合；文中也检验热流和离子黏性应力的适用性，不能仅因有激波解就默认流体闭合在任意薄层都成立。
\loc{L164}{PDF第10--12页式(4)--(18)，第16--20页图4、5及物理解释；后续图6、7检查热流/应力，图9--11给非稳态过程}

\subsection{电场梯度怎样破坏离子的漂移近似}
Yoon与Bellan（2018，\texttt{Yoon2018Heating}；2021，\texttt{L272}）强调的是\emph{电场空间变化相对回旋响应有多快}，不是只比较电场幅值。给定近均匀$\B=B_0\hat z$，低频垂直电场的首阶响应为$\vct v_E=\E\times\B/B^2$，下一阶极化漂移为$\vct v_p=m_i\dd\E_\perp/(q_iB^2\dd t)$。梯度很强时，沿轨道看到的场变化使这个“下一阶”不再小，离子无法保持稳定回旋中心和磁矩。局地量级指标为
\begin{equation}
\alpha_i=\frac{m_i|\nabla_\perp\cdot\E_\perp|}{|q_i|B^2},\qquad
\Omega_{ci}=\frac{|q_i|B}{m_i}.
\end{equation}
它无量纲；增大质量电荷比使失磁化更容易，增大局地$B$倾向于恢复磁化。\textbf{但绝对值指标是特定几何的诊断，不是对任意电场保证正熵产的充分判据。}例如只取一维局部$E_x=Gx$，相邻轨道满足
\begin{equation}
\delta\dddot x=\left(\frac{q_iG}{m_i}-\Omega_{ci}^2\right)\delta\dot x.
\end{equation}
此式中梯度的\emph{符号}决定局地指数不稳定性，不能用$|G|$抹掉恢复力与发散力的区别；无界轨道指数分离也不是有界混沌或热化的完整证明。这是本报告从运动方程作的局部教学推导，须与原文数值轨道证据分开。\citep{Yoon2018Heating,L272}

2018年的电子流体场呈入流向势谷、出流向势丘；离子跌入势谷后沿出流势降增能，增长场使初始相邻离子的经历不同。图1是势结构，图2是轨道，图3是流出速度展宽，三者不能用一张温度图替代。固定离子EMHD只生成背景场，测试离子无反馈，故不能同时闭合离子与电子能量。2019年\texttt{F743}的PIC图1--5进一步对照$\alpha$、温度和粒子轨道，并扫描$B_g/B_0=0$至0.5；较强导场抑制该机制，但残余平行温升区也提示存在其他通道。该趋势不能外推为所有强导场装置的离子温度都必须降低。\citep{F743}
\loc{Yoon2018Heating / L272 / F743}{依次图1--3及式(2)--(17)；正式版PDF第3页式(1)--(13)、图1；2019图1--5及静态Harris阈值的适用性讨论}

\subsection{补上的拾取参照：Drake2009与Yoon2023并非完全同一解释}
Drake等（2009，\texttt{Drake2009Pickup}）的反平行PIC显示：离子跨过狭窄出流边界时，磁场被Hall面外场迅速转向，粒子尚未来得及跟随$\E\times\B$速度。进入快速运动的磁通管后，它相对于出流具有大速度，表现为拾取状回旋/反向束。图2把轨道、磁矩破坏和分布放在一起，图3变换到随出流运动的de Hoffmann--Teller参考系，使局地电场近零；在该系中能量近守恒而运动方向改变，变回X线系后得到体流加速和中心矩展宽。\citep{Drake2009Pickup}

冷入流、近稳态、弱导场近似给出原文式(7)、(8)：
\begin{equation}
\Delta T_\parallel\simeq m_i v_0^2\frac{B_{0L}^2}{B_{0L}^2+B_{0M}^2},\quad
\Delta T_\perp\simeq\frac{m_i v_0^2}{2}\frac{B_{0M}^2}{B_{0L}^2+B_{0M}^2},\quad
\Delta T_i\simeq\frac{m_i v_0^2}{3}.
\end{equation}
温度以能量计，$v_0$为出流速度，$B_{0M}$是下游Hall场，不是可任意替换的外加导场。这是二阶矩预测，不要求Maxwell分布；文中PIC质量比25，且预测的各向异性并未精确吻合。附录A的“反弹频率”定义为$v_0/\Delta$，仅在指定极限才是实际频率，不能当成一般精确周期。

\textbf{必须保留的研究分歧。}Drake2009第16段在下游一维近稳态模型中认为，面内静电场在往返轨道上增能后又回收，面外重联电场对净位移做功。Yoon等（2023，\texttt{Yoon2023PhaseSpace}）第5节则批评把拾取作为零/弱导场重联的总解释，认为其模拟只有初始短暂拾取状阶段，随后主要是单阱中的轨道混合；作者还在结论中将净增能归于面内场，称重联电场净功为零。\emph{这是原文之间的归因差异，不能被报告写成已统一的事实。}2023图6的四阶段与图7附录的分力速度积分是具体证据，但速度分量末值恢复并不单独推出该电场的累计功为零。可检验的共同诊断应是对同一粒子、同一时间窗计算$\int q_i E_Nv_N\dd t$、$\int q_i E_Lv_L\dd t$、$\int q_i E_Mv_M\dd t$，并核对总动能增量；磁力只改变方向，分量间能量转移相消。本报告未复跑两文轨道，不能裁决其全部差异。\citep{Yoon2023PhaseSpace}
\loc{Drake2009Pickup / Yoon2023PhaseSpace}{前者全文14页含附录A；图1--4、式(7)--(10)、第16及31--39段；后者PDF第6--8页图6、7及第5节，亦见图3的冷多束}

\subsection{电子为何在重联以后、沿分离线才热化}
Zocco与Schekochihin（2011，\texttt{E512}）及Loureiro等（2013，\texttt{Loureiro2013Heating}）的强导场图像是：重联先建立流动和扰动，沿磁力线速度不同的电子逐渐错相；同一空间位置的分布在$v_\parallel$方向出现越来越细的纹理。用Hermite阶数$m$表示这些速度结构时，能量从低阶的流速、温度矩流向高阶；有限碰撞对细速度结构作用更强，才将这份非平衡自由能转为背景热。\citep{E512,Loureiro2013Heating}

以原文自由能的体积平均约定记
\begin{equation}
H_e=\frac1V\int\dd^3x\dd^3v\frac{T_{0e}g_e^2}{2F_{0e}},\qquad
\frac{\dd W_{\rm fluid}}{\dd t}=-\mathcal Q-D_\eta,\qquad
\frac{\dd H_e}{\dd t}=\mathcal Q-D_{\rm coll}.
\end{equation}
这里$F_{0e}$是背景Maxwell分布，$g_e$是除去低阶流体响应后的扰动；$\mathcal Q$是压力功导致的可逆交换（可变号），$D_\eta$是电阻耗散，$D_{\rm coll}=-V^{-1}\int T_{0e}g_e C[g_e]/F_{0e}\dd^3x\dd^3v\geq0$由满足H定理的碰撞算子给出。严格零碰撞时，$H_e$增加仍不意味着精细熵增加。\textbf{这里写的是扰动自由能预算，不能与第2章完整物种内能重复相加。}

Loureiro2013图2左列叠加磁通等值线与电流，中列叠加电势与电子平行温度扰动，右列为耗散；上、下两行分别在重联峰$22.7\tau_A$与耗散峰$29.4\tau_A$，应横向、纵向都比较。图4底行显示速度细纹形成，上两行显示自由能到高$m$再耗散，因而“先重联、后热化”具有可见证据。Tassi等（2018，\texttt{L244}）有限矩热流把这条链压缩为即时闭合，图4却缺少同样的峰时滞后，原文将是否过耗散列为问题。Numata与Loureiro（2015，\texttt{Numata2015Heating}）恢复双物种GK后，高$\beta$下离子平行/垂直相混更重要；电子多在分离线/岛缘、离子可在次级磁岛内部耗散。因此“电子用Landau、离子用Hall”不是完整分类：空间结构、相空间通道和排序共同决定可见结果。\citep{L244,Numata2015Heating}
\loc{E512 / Loureiro2013Heating / L244 / Numata2015Heating}{依次PDF第11--12页式(76)--(84)；图1、2、4、5；图4及结论；式(2.5)、图5--7与11--12、附录A}


\chapter{实验现象究竟验证了什么}
\section{MRX：电势加速和下游热化不是同一地点}
Yoo2013零外加导场氦实验中，面内电场约500--800 V/m，较重联电场约200 V/m更大；电势差约10--35 V，使离子在面内加速至约半个上游Alfvén速度。主要离子温升出现在更下游、流速趋饱和处。电势由浮动电势和电子温度换算，另由电子动量平衡积分作一致性检验；流速图主要用光谱标定的Mach探针，温度由精细结构拟合后的光谱线宽得到，诊断细节见第\ref{sec:measurement_chain}节。\citep{Yoo2013}

Yoo2014的F679增加了实验细节及匹配碰撞量级的VPIC对照。无离子碰撞时，下游二阶矩含多束；加入现实量级碰撞后更接近Maxwell。再磁化、势阱及停留时间使碰撞有效。它支持“势能加速---轨道混合---下游碰撞热化”的条件性链条；不证明所有无碰撞Hall电场都立即产生熵。前后论文继承同一装置/数据线，不算独立装置重复验证。\citep{F679}

电子经典电阻估计约占观测沉积的17\%，余量提示缺失过程，但压力张量、热通量、涨落和损失未全部直接闭合。不能仅凭经典预算不足就唯一宣布某一种异常电阻机制获证。
\loc{Yoo2013 / F679}{前者PDF第2--4页图2--4；后者第5、8--11页图5、9--13及式(1)--(19)}

\section{MRX常见的2:1应怎样引用}
Yamada等2015正式版F674图7的功率预算为下表，单位MW；原文术语的定义一并保留。\citep{F674}
\begin{center}\begin{tabular}{P{74mm}rr}
\toprule 量 & 离子 & 电子\\\midrule
电场能量沉积率$W_s$ & $0.71\pm0.15$ & $0.35\pm0.06$\\
有序流动能变化 & $0.14\pm0.06$ & 近零\\
文中“热能变化”$\Delta W_{H,s}$ & $0.27\pm0.10$ & $0.26\pm0.10$\\
估计热损失 & $\leq0.35$ & $0.15\pm0.08$\\\bottomrule
\end{tabular}\end{center}
所以约2:1对应沉积率，而热能变化约1:1。原文式(6)的$\Delta W_H$\textbf{还含$3nT/2$储存及$5nT\vct u/2$边界焓流差}，不能等同于本专题$\Delta\mathcal U$，更不能等同$Q^{\rm irr}$。误差与损失上限使图内分量不是精确代数恒等式，应报告残差而不是强行补齐。

表I以输入磁能为分母：MRX控制体从$1.5d_i$变为$2d_i$，离子/电子沉积份额由约0.31/0.21变为0.37/0.18；匹配MRX的PIC约0.26/0.19，开放PIC约0.29/0.11。控制体、边界、装置和初温比不同，比例随之改变。F674正式稿与旧机构预稿同属一篇；其2014年前序研究也不能作为第二组完全独立数据。
\loc{F674}{11页正式版，第6--9页图6--8、表I及式(2)--(6)}

\section{非对称MRX进一步限制普适比例}
Yoo2017密度比约8的实验中，低密度侧分离面附近电子升温更强。图10给出离子沉积$0.44\pm0.09$ MW、电子$0.35\pm0.07$ MW，两物种$\Delta W_H$均约$0.30\pm0.06$ MW。电场沉积相近，与对称案例的2:1不同。\citep{Yoo2017Heating}

二维和三维VPIC的温度增量相似，但三维分布更接近Maxwell；LHDI可改变热化程度而不显著增加净温度。实验波动与温升相关系数约0.29、$p\simeq0.086$，不能据此唯一确立波动因果机制。这恰好说明“增能多少”和“分布怎样变热”是两个验证目标。
\loc{Yoo2017Heating}{18页全文；图4、5、8、10、11与表1；PDF第15页预算图10（原笔记页码已纠正）}

\section{MAST：强导场下分别测到热点和交换时间}
Tanabe等2015通过Thomson散射和离子光谱层析得到：X点电子热点宽约2--5 cm，电子温度可在小于1 ms由约10升至200 eV；离子高温区更宽，位于出流及停滞区；经过约4--10 ms物种能量交换，两种温度可呈三峰。论文直接引用L223的双流体模拟，支持下游黏性/压缩和热输运解释的空间合理性。\citep{Tanabe2015MAST,L223}

但二维剖面组合多次脉冲，离子诊断用杂质谱线；缺乏直接磁探针使局地机制不能唯一确定。导场0.3--0.8 T、重联场约0.07--0.15 T，强导场与闭合磁面改变热输运及驻留时间；在所测条件下更强导场使电子峰更尖，下游离子温度基本不变。它与弱导场PIC中梯度随机增能被抑制的趋势不构成直接矛盾，因为比较的场景、量和时间窗不同。
\loc{Tanabe2015MAST}{5页正式稿全部，PDF第2--4页图1--5，参考文献[41]直接连接L223}

\section{HIT-SI、MMS与TS-6：分清验证等级}
HIT-SI已有三维双温Hall与多诊断的比较，主要约束整体温度、流动、密度和驱动响应；尚非全面局地熵产测量。MMS事件的十矩比较则有温度各向异性幅度合理、空间侧别不符的明确成败对照。这两种结果都比仅展示四极磁场更接近加热问题，但仍不足以唯一标定所有闭合。\citep{Kaptanoglu2020,TenBarge2019MMS}

本轮也阅读了2025年IAEA会议TS-6稿。引言提到Hall MHD和PIC，实际模型/结果段给出的是VPIC，未展示独立Hall方程、网格和加热结果；所谓约20\%符合针对PIC电势幅度，而非两物种热化比例。文中$\Delta T_i$形式未明确扣除平均流且缺少三维中心矩的常规因子，不能直接当作本报告定义的热温度。会议稿可作为待后续完整验证的线索，不能用其题名与摘要补足缺少的Hall结果。\citep{TS6FEC2025}
\loc{TS6FEC2025}{8页会议稿全文，模型和结果第4--7页；尚未取得相应独立Hall计算或源数据}

\section{按诊断、图像和推断层级重读实验}
\label{sec:measurement_chain}
本节补足前述结果的观测链。“测量”也经过仪器响应和拟合；这里将\textbf{仪器信号/直接约束、派生量、机制解释}分别说明，避免把反演图当成逐点直接探测。除原文报告的误差外，不自行补造精度；以下分辨率不是所有测量量共有的误差条。

\subsection{Yoo2013/F679：按四张图还原MRX离子的经历}
\textbf{装置和诊断。}Yoo等（2013，\texttt{Yoo2013}；2014，\texttt{F679}）采用4.5 mTorr氦放电、无外加导场、磁通芯间距42 cm；后者给$d_i\simeq8$ cm。磁探针取得$\B$，其空间导数给$\J$，磁通的时间导数给面外重联电场；Langmuir探针给$n_e,T_e$。浮动电势$\Phi_f$经
\begin{equation}
\Phi_p\simeq\Phi_f+[3.3+0.5\ln(m_i/m_p)]T_e
\end{equation}
换算等离子体电势（此处$T_e$数值用eV、电势用V），再由$-\nabla\Phi_p$得到面内场。\textbf{电子动量平衡不是这一电势图的唯一来源}：作者另由$\vct u_e=\vct u_i-\J/(en)$、磁场和电子压强估计电场，积分后与探针换算电势比较；Yoo2013图2和F679图5的点线对照正是这个检验。\citep{Yoo2013,F679}

离子谱线He II 468.6 nm的中心位移约束视线流速，宽度约束温度；必须拟合13个精细结构分量，否则温度高估15--25\%。实际展示的较细流速图以经光谱标定的Mach探针为主。光谱空间平均约3--4 cm、曝光5.6 $\mu$s；磁探针最细径向0.6 cm、轴向3 cm。F679从4200余次放电筛出约30\%，每点平均7--15次，误差取样本标准差与诊断不确定度中较大者。因此二维图是重复放电拼接，有可重复性筛选，不是同一次放电的瞬时三维快照。

\textbf{读图顺序。}先读Yoo2013图2：势阱横向束缚、向下游变深变宽；再读图3(a)：离子穿越分离线后转向出流，图3(b)显示面内电功集中在下游。最后读图4：下游谱线既偏移又展宽；沿$Z$的流速先升、约$Z>9$ cm趋饱和时$T_i$才明显上升（图示由约数eV到约15--20 eV，非精确拟合数值）。温度横向峰比出流峰窄，因而不能只用$u_i^2$解释温度图。这是加速与温升空间分离的直接约束。

\textbf{哪些解释仍有竞争。}Yoo2013明确指出，所测热点不在经典无磁化黏性剪切最强处，经典热传导又难维持该梯度，故不能把这份MRX温升简单贴上“经典黏性加热”。原文提出再磁化延长驻留、束与背景碰撞、波粒散射、势阱反弹和磁通芯边界等可能性；上游离子平均自由程6--12 cm、下游2--5 cm支持下游碰撞有效。F679图9--10再以有/无现实量级离子碰撞的VPIC对照：两者均可有较大二阶矩，有碰撞时分布更接近Maxwell。\textbf{这份分布对照来自模拟；实验光谱的一维近Maxwell形状并不等于直接测得完整三维分布和熵。}同装置前后工作有数据继承，不能当作独立装置重复。
\loc{Yoo2013 / F679}{前者PDF第2--4页图2--4；后者第3--5页诊断与电势、第8--9页图9、10；第4页误差选取标准}

\subsection{MRX电子：电功集中处、温度峰和热量去向分开读}
F679图11显示电子由上游5--7 eV到分离线内11--12 eV，温度峰略在电子层外；图12的$\J_e\cdot\E$则集中于X点邻域，原文估计约70\%的电子沉积在小蓝框内。沿场传导把能量带离局地，所以温度最大处并非电功最大处。式(8)给$340\pm30$ kW电功，式(10)经典电阻约$59\pm12$ kW。约17\%的不足说明经典模型未闭合所测预算，却不能把剩余83\%直接分配给某一种波或“Hall电子热源”。\citep{F679}

\texttt{F674}（Yamada等，2015）把这种局地信息组成有限体积功率账本：$31.5\leq R\leq43.5$ cm、$0\leq Z\leq15$ cm，按环向对称积分；准稳态仍保留储存项，压力近似各向同性。局地关键量约10--15\%误差；热流采用Braginskii估计，氧杂质辐射由谱线与浓度估计，电荷交换损失因未测中性密度而取上界。图7数值因此混有测量、导数/积分、输运模型及上界，绝不是四个独立热量计读数。用输入Poynting功率作分母与用净转换磁功作分母也不同，且出流Hall场带走的Poynting通量不可丢弃。\citep{F674}
\loc{F679 / F674}{前者PDF第9--11页图11、12及式(6)--(19)；后者正式PDF第7页式(2)--(6)与损失估计，第8页图7，第9页表I}

\subsection{Yoo2017：大电场不一定意味着该处离子电功大}
Yoo等（2017，\texttt{Yoo2017Heating}）控制密度不对称：典型高/低密度侧约$4/0.5\times10^{13}\,\mathrm{cm^{-3}}$、重联场145/168 G，无外加导场。约2200次放电筛出35\%，每点7--10次，典型标准差10--15\%；离子光谱曝光10 $\mu$s、空间平均4 cm。图2显示强势垒偏向低密度分离线，使高密度侧来的多数离子留在出流内；图11因此给出一个反直觉但清楚的现象：\textbf{低密度侧面内场虽可约600 V/m，离子流向与电场未有效对齐，该处离子净电功仍小。}高密度侧场约150 V/m，离子电功约15 W/cm$^3$，也低于对称参照。电功需要$\J_i$与$\E$的点积，不能按场强排名。\citep{Yoo2017Heating}

图4电子热点却向低密度侧分离线延伸。作者逐项估计极化、梯度和曲率漂移，在该观测位置不足以解释温升，提出$E_\parallel$加速与有限碰撞的解释；$E_\parallel\sim40\pm20$ V/m是从较大总场投影的小量，误差不能省略。约4 cm平均自由程内的能量增量约2 eV，与观测量级相容。这是位置限定的机制估计，不能推广为所有区域曲率/Fermi通道都无效。

图8--9中的\emph{总电子电功}主要还含抗磁电流与垂直电场的功。原文式(7)--(10)用压力平衡、焓输运解释它，不能将其全部等同于上段单粒子的平行增能。本报告采用完整内能方程区分压力功和通量，亦不把一般压力功一概宣布为不可逆或“完全不加热”。图5的2D/3D VPIC给近似温度增量、不同Maxwell化程度；图6实验波动与温升弱相关（$r\simeq0.29,p\simeq0.086$），支持保留LHDI的可能作用，但未确立唯一因果。图10两物种功率更接近，既与电子低密度侧附加沉积有关，也与离子跨势垒路径减少有关，不能只说“非对称性使电子更热”。
\loc{Yoo2017Heating}{PDF第4页诊断，第5页图2，第7--10页图4--6及机制估计，第12--14页图8、9及式(7)--(10)，第15页图10、11；图10不是PDF第12页}

\subsection{MAST：先建立两个不同热点，再观察能量交换}
Tanabe等（2015，\texttt{Tanabe2015MAST}）图1中P3线圈驱动两环沿竖直方向合并，径向喷流进入下游闭合磁面。图2是单次放电8帧Thomson序列，0.1 ms间隔，电子温度从约10到200 eV发生在不到1 ms内，并与Mirnov磁信号相联系。图3的二维图则由多次放电移动采样位置拼接：电子窄峰宽2--5 cm，离子热点在下游更宽；不能把两类时间/空间成像方法混成单次全场电影。\citep{Tanabe2015MAST}

电子诊断为300通道Ruby和130通道YAG Thomson散射，含光学模糊的空间分辨约15/10 mm；离子为32弦Doppler层析，主要CVI 529.05 nm，早期冷帧用CIII。离子图的温度是杂质谱线反演，向主离子的解释需依赖碰撞与诊断适用性。图3离子热点颜色标度约百eV量级，不能拿L223高驱动模拟的keV峰冒充这一实验的实测值。图4给出离子下游双峰与电子中心峰在约4--10 ms交换时间后趋向三峰，图5在近固定$B_{rec}\simeq0.08$ T下改变导场0.4--0.8 T，电子中心峰更尖、所测下游离子温度基本不变。

作者给出两种使电子峰变尖的可能解释：更强导场降低跨场热输运、延长高电场区驻留，或减少meandering幅度使电流层更尖。缺乏直接磁探针，本文不能在两者之间作唯一判决。$B_{rec}$又由合并后30 ms的EFIT估计，并非每个加热时刻局地测量。对离子，强导场虽可抑制某些黏性阻尼，但闭合磁面延长驻留，仍可能最终耗散出流。\textbf{实验约束的是峰位、宽度、时序和参数趋势；没有直接给出局地黏性熵产、完整各向异性张量或普适两物种热化比例。}
\loc{Tanabe2015MAST}{PDF第2页诊断与几何，第3页图2--4，第4页图5及竞争解释；第5页原文引用[40]、[41]连接Stanier和Browning模拟}


\subsection{HIT-SI：高频驱动下的热离子与冷电子是哪一层证据}
Kaptanoglu等（2020，\texttt{Kaptanoglu2020}）第III.A节引用HIT-SI高频68.5 kHz测量，$T_e\sim5$--10 eV、$T_i\sim20$--30 eV；并提及\emph{另一个装置HIT-SI3}在14.5 kHz的初步Thomson电子温度5--10 eV。不能将后者写成HIT-SI同一次放电的共址双物种测量。作者主要比较光谱、干涉、磁探针的整体趋势；原始诊断学位论文及HIT-SI3会议结果本轮未独立核读，精确温度数值仅是该模型论文转述的量级约束。\citep{Kaptanoglu2020}

模型图2的密度是合成干涉弦平均；温度与图3功率大多是\emph{包含注入器}的体积量，图4才给空间切面。随着驱动频率增加，模型离子温度明显增加、电子较弱变化：场反转更频繁，黏性与压缩使离子增热，电子欧姆热增加的同时壁面热流也增加。温度是源和损失平衡，不能把低$T_e$解释为电子从未获得能量。装置持续1--2 ms，物种交换约100--200 $\mu$s，仍可慢于注入周期，支持独立两温方程。作者明确因缺少真实注入电路而避免严格定量验证；勘误正文也未取得，所以这里不从上述转述数值提炼精确加热比。
\loc{Kaptanoglu2020}{35页接收稿第5页式(16)、(17)与边界，第6--8页第III.A节及图2，第9页图3、4；第15页关于缺电路而不作定量验证的说明}

\subsection{新增空间观测参照：同一标度并不保证系数正确}
Drake等（2009，\texttt{Drake2009Pickup}）图6用Wind的22次磁剪切角大于$120^\circ$的出流穿越检验$\Delta T_p\propto(\Delta v_L)^2$；相对于式(8)的预测，拟合斜率只有0.39，不能只引用标度“吻合”而隐去系数偏低。图5示例的质子/氦核温升约5.5/20 eV，而预测约14.7/58.6 eV。图7及表I的ACE重离子用SWICS 12 min分辨率的径向$T_{rr}$，不是完整标量温度；第一次短穿越无居中采样，后两次相对局地磁场方向也不同。\citep{Drake2009Pickup}

这些是空间观测而非受控实验。作者把远下游边界结构、各向异性散射及膨胀列为解释缺口；本报告据此要求流体闭合同时检验总增量、各向异性和采样位置。此文不提供电子分配，也不能填补强导场2009年ApJ论文尚未核读的缺口。
\loc{Drake2009Pickup}{PDF第8--10页图5--7、表I，第11--12页讨论；SWICS只测径向温度的限制见第9页}


\chapter{已经解决什么，为什么还有问题}
\section{报告缺口与研究未决问题分开登记}
\label{sec:gap_map}
本节是2026年11月13日的实际复查结论。\textbf{“报告缺口”意味着现有原文已有证据而旧正文未讲清；“研究未决”意味着所读原文仍未建立该结论。}前者本轮已修订，后者必须保留；新检索未发现完整证据也不等于全领域无人解决。除下列新增Drake2009全文外，本轮对旧文是针对正文、方法与关键图页的复查，不将其重复计为新增全文核读。
\begingroup\small
\begin{longtable}{P{30mm}P{59mm}P{59mm}}
\toprule 问题与证据 & 旧报告的不足、本轮已补 & 仍存边界及判断归属\\\midrule\endfirsthead
\toprule 问题与证据 & 旧报告的不足、本轮已补 & 仍存边界及判断归属\\\midrule\endhead
双温Hall热点；Browning2016（L223）图2、3 & 补磁张力出流、离子黏性、电子超电阻及热输运因果链；不再把电子峰概括为简单欧姆热点 & 【原文】初步双温、超电阻影响强。【综合判断】峰位相容未唯一验证耗散闭合。\\
离子梯度机制；Yoon2018、2021（L272）、F743 & 补势谷/势丘、漂移层级、$\alpha$与带符号局部不稳定性；分清轨道分离和温度 & 【原文】不同导场/区域仍有其他通道。【综合判断】$\alpha>1$、轨道混沌与熵产不能直接画等号。\\
拾取与轨道混合；Drake2009Pickup、Yoon2023\allowbreak PhaseSpace & 新核读14页含附录；补二阶矩公式、功归因分歧与共同轨道积分检验 & 【原文】2009标度系数不足；2023机制份额未定。【综合判断】需同轨道预算，当前未裁决分歧。\\
电子相混；E512式(76)--(84)、Loureiro2013图2、4 & 补流体自由能到高阶矩再碰撞耗散；解释热化延迟与分离线位置 & 【原文】L244质疑无延迟与过耗散；Numata保留速度收敛限制。【综合判断】不能由相近重联率验证热化时序。\\
MRX诊断；Yoo2013图2--4、F679第II节 & 纠正电势来源；补谱线拟合、分辨率、放电筛选和温升先后；明确经典黏性解释的局限 & 【原文】热化有竞争过程，需模拟辅助。【综合判断】谱线一维拟合未直接测完整熵；与MAST黏性解释不可混用。\\
MRX分配；F674图7、Yoo2017图10 & 补控制体、损失反演、误差和大电场但小离子功的原因；纠正图10物理页码 & 【原文】中性损失取上界、热流依模型。【综合判断】实验功率比与不可逆热化比仍不同。\\
MAST温升；Tanabe2015图1--5 & 补单次时间帧与跨脉冲二维图差别、杂质诊断、EFIT时刻；串联双峰到三峰 & 【原文】缺直接磁探针、峰尖化有竞争解释。【综合判断】驻留与热输运可能掩盖局地源项变化。\\
HIT-SI与MMS；Kaptanoglu2020、TenBarge2019MMS & 保留装置几何/驱动损失及MMS各向异性侧别成败，明确是局部验证 & 【原文】电路/边界与闭合限制。【文献缺口】勘误正文、正式版新增图未核读；不能用精确比例作强结论。\\
\bottomrule\end{longtable}\endgroup

\section{从理论到观测：用什么现象区分解释}
\begin{longtable}{P{35mm}P{55mm}P{58mm}}
\toprule 物理图像 & 已有现象与原文定位 & 还需什么才能加强判定（本报告建议）\\\midrule\endfirsthead
\toprule 物理图像 & 已有现象与原文定位 & 还需什么才能加强判定（本报告建议）\\\midrule\endhead
电势先加速，后混合/碰撞 & MRX流速先升、下游谱线才明显展宽；Yoo2013图4、F679图10 & 同时看平均流、完整分布与停留时间；只拟合$T_i$无法区分冷多束和真正分布展宽。\\
出流黏性/压缩、物种交换 & L164图4的三段结构；MAST图3、4不同初始热点、晚时三峰 & 对齐几何、热边界和交换时间，测速度梯度及热流；压缩升温中可逆部分另列。\\
电子相混后弱碰撞耗散 & Loureiro2013图1、2、4的延迟和高阶谱；L244图4没有相同延迟 & 这是理论/数值对照，尚非所列装置直接观测电子Hermite级联；需分布和耗散时序验证。\\
压力张量与热流闭合 & TenBarge2019图1有物种内能；图6--10各向异性量级合理、侧别错误 & 同时匹配侧别、峰宽、热流和能量守恒，不能凭一个峰值宣布十矩已充分验证。\\
局地电功与非局地输运 & F679图11、12电功与温度峰不同；Yoo2017图8、9抗磁功大 & 联合压力功与焓/热流预算；正$\J_s\cdot\E$仅证明场向粒子输能，不证明局地不可逆热化。\\
\bottomrule\end{longtable}

上述对应也显示覆盖的真实不足：本专题对MRX/MAST的联系最具体，HIT-SI是装置整体模型比较，MMS主要依TenBarge文的事件对照；并未独立全文核读所有原始MMS观测、所有诊断方法论文及其他装置的加热系列。强导场拾取阈值、Fermi/俘获等更广的动理学加热综述仍不能由这一专题替代。它们应列为定向追读，而不是用未读摘要填进“已验证机制”。

\section{在明确条件下已经建立的结果}
\begin{longtable}{P{34mm}P{62mm}P{52mm}}
\toprule 问题 & 已达到的程度与证据 & 不能顺带推出什么\\\midrule\endfirsthead
\toprule 问题 & 已达到的程度与证据 & 不能顺带推出什么\\\midrule\endhead
流体能否分别计算温度 & L223双温Hall局地热点、Kaptanoglu2020三维装置计算；TenBarge2019分物种热能预算 & 不能因能解两温便认定闭合已普遍验证。\\
离子怎样先加速再变热 & Yoo2013/F679电势、流速和温度空间关系，加碰撞VPIC出现Maxwell化 & 尚非所有环境唯一热化机制，更非裸Hall熵产。\\
为何离子/电子热点不同 & L164碰撞激波/传导的时序；L223与Tanabe2015的局地电子/下游离子热点 & 热输运、边界与驻留时间不同，不能只比较峰值比。\\
电子相混能否通向不可逆热 & E512的自由能恒等式；Loureiro2013的非线性Hermite和延迟耗散 & 所选低$\beta$模型没有完整离子竞争通道。\\
有限矩是否有可改进空间 & L244引入Landau耗散；TenBarge、Ng2020、Closure2026展示闭合改变模态/热结构 & 更多矩或相近重联率不保证熵产及非热分布正确。\\
实验能否给物种预算 & F674/F679/Yoo2017在有限控制体给沉积与热能/输运估计 & 不可把2:1或近1:1混写为普适不可逆热化比。\\
高二阶矩是否总意味着热化 & Yoon2023显式区分冷多束和混沌展宽；F679与Yoo2017比较分布 & 不能把任意压力张量峰或正电功图称为热熵已增加。\\
\bottomrule\end{longtable}
表中结论均有上述章节的原文定位和引文；其中条件性“已解决”是建立机制、预算或可计算模型，不代表消除了所有不确定性。

\section{未解决项按原因分类}
\subsection{模型所保留的信息不足}
若只有等温电子或固定$T_i/T_e$，温比就是输入，问题在所选模型定义；提高网格不能把不存在的未知量算出来。有限矩不能唯一恢复多束、尾部和速度空间细化；E512/L244主动删去的离子热化通道、EMHD固定离子和试粒子无反作用也属范围限制。需要增加群体、矩或动理学自由度，且代价与新闭合随之增加。F635改进多群表达仍不含自洽离子热化，是具体例子。

\subsection{闭合的可迁移性未建立}
异常电阻、超电阻、黏性和热流若为调参，得到的温度比可能主要反映这些输入。L223的电子局地耗散受超电阻控制，TenBarge2019的固定尺度闭合压制LHDI，L244缺少电子热化延迟，都是可定位的例子。原文/结果提出了问题，本报告综合判断是：需在多个几何、导场、$\beta$和碰撞率下进行\emph{不重新拟合}的外推验证，而非只重现一张温度图。

\subsection{数值误差、尺度和耗散会污染分配}
总能相对初始值误差小，不代表相对\emph{释放的磁能}误差小。网格耗散回不回内能、回到哪个物种，是能量守恒与热化解释的关键。五矩MMS运行、L244的释放能比例误差、有限质量比和未解析回旋半径，以及Hermite/速度网格截断，均可能改变通道。数值收敛只验证离散逼近了选定方程，不能单独证明闭合符合真实等离子体。

\subsection{实验比较还有定义和可观测量限制}
电势反演、视线积分、杂质温度、跨脉冲拼图及边界热损失不确定性，使局地预算不完全可测。MRX对称/非对称比例差异、MAST强导场驻留和MMS侧别错误说明必须对齐密度、初温、导场、几何及仪器平均。缺少完整压力张量和热流时，匹配温度幅度不能唯一认证机制。

\subsection{有些问题只是尚未在所读论文中检验}
Finelli2021明确将能量分配留待后续；Yoon2023留下精确能量份额与机制份额；L247验证算例强制单温。它们不是证明对应模型原则上不可能，而是没有提供用户所需的验证。勘误或正式版未取得是文献证据缺口，也不应伪装成物理反例。

\section{对“能不能研究加热比例”的最终判据}
若目标是\textbf{给定装置、给定输运模型下的温度和功率预算}，独立两温Hall/双流体有明确用途；应把它当成可检验的条件性预测。若目标是\textbf{无碰撞重联中无需调参的不可逆离子/电子分配}，现有证据要求更多动理学信息或经过独立验证的闭合，标准Hall项本身远远不够。若目标是\textbf{非热尾、束与热化路径}，单一标量温度甚至十矩也可能不够，需要直接诊断分布或相应扩展。

不是所有未决问题都出在Hall/双流体方程上：模型信息、闭合、本构参数、数值误差、实验诊断、边界输运和未完成验证必须分别定位。把这些差异压成一个“Hall可以/不可以”会掩盖已经取得的进展，也会夸大尚未取得的能力。

\chapter{怎样开展一项可判定成败的加热研究}
以下为本报告提出的研究设计，\textbf{不是已经运行的数值实验}。先选择一个实验几何或公开动理学基准，明确研究的是局地温升、控制体功率，还是熵产比例；避免一开始寻找跨所有场景的常数。

\section{最小可解释的双温Hall实现}
演化密度、离子动量、磁场及独立$U_i,U_e$；用$\vct u_e=\vct u_i-\J/(en)$处理电子对流与热流，不能直接令两种速度相同。由所选广义欧姆定律更新磁场；电阻/惯性/压力项和内能源要在离散层面成对一致。标量模型中$U_i,U_e$各是一条独立热力学方程；十矩模型中改为演化压力张量并用迹计算内能。

明确写出$\eta,\nu_i,\nu_e,\kappa_\parallel,\kappa_\perp$、物种交换率、辐射/电离和壁面损失是否存在，来源是经典碰撞理论、测量、拟合还是数值正则化。记录磁场/速度离散位置、散度约束、Hall刚性时间尺度、显隐式分裂、边界热流及能量回收步骤。L247可作输运实现参考，但它的单温验证不足以替代新增双温验证。

\section{必须一起输出的诊断}
\begin{enumerate}
\item 两物种$n,\vct u,\vct P,T$及磁场；若有分布，另给多束分解、各向异性、非热尾和粗粒化定义。
\item 分别输出$\J_i\cdot\E$、$\J_e\cdot\E$、$-p_s\nabla\cdot\vct u_s$、$-\vct\Pi_s:\nabla\vct u_s$、碰撞交换、物理耗散、热流及边界焓流。
\item 每种粒子流能与内能、电磁能、输入和输出；残差同时用初始总能及释放磁能归一，避免小分母被遮盖。
\item 同一固定控制体、同一时间窗的$R_W,R_U,R_Q,R_T$；其中无法可靠定义或测出的量直接标缺失，不用另一比值代替。
\item 若使用Landau/高阶矩，分开记录流体自由能、高阶矩自由能、碰撞或数值截断耗散；若使用压力松弛，核对零迹及正性。
\end{enumerate}

\section{验证顺序与可证伪结果}
先做选定离散能量账本、均匀物种温差交换和物理耗散的守恒检查，再改变空间/时间分辨率、人工耗散、质量比和边界距离。不要只因重联率稳定就停止热化收敛检查。

随后在相同初边值比较双温Hall、十矩及动理学：首先看能量预算是否闭合，再看热点位置/宽度、两物种随时间的增量和分布，最后看熵产。闭合参数一旦在训练案例选定，应固定后用于其他导场、$\beta$、碰撞率和不对称性；若必须逐例调参，应报告为校准模型而不是通用预测。

实验对照需生成相同视线积分、采样位置、时间平均和误差条。针对MRX，要检验加速区与热化区的分离以及有/无碰撞分布；针对MAST，要检验电子窄峰、离子下游峰和交换时间；针对MMS，要同时检验温度幅度与所在一侧；针对HIT-SI，要纳入驱动、壁面和多诊断约束。

如果双温Hall在能量守恒且网格收敛后仍系统性漏掉分布多束或热点侧别，才可将差距进一步归于闭合/模型信息；若控制体热流或损失仍未测定，不能提前归咎于Hall项。这一研究顺序能把“算出了不同温度”推进到“有证据支持不同物种加热的解释”。

\appendix
\chapter{证据版本、逐篇笔记和维护入口}
本附录按正文首次引用排序。所有引文均已列入全库核读档案；本轮未重新通读的旧证据明确标识。点击本地PDF需保持完整目录结构。详细阅读笔记含方法、参数、图式、限制及未取得材料。

\paragraph{L223（2016）}\label{paper:L223}
Two-fluid and magnetohydrodynamic modelling of magnetic reconnection in the MAST spherical tokamak and the solar corona。11月12日新增全文核读；本次修订不重复计数；主版本16个PDF物理页。
arXiv:1507.07432v1 (2015-07-27); article published 2016。\href{run:../../papers/L223.pdf}{打开本地原文}。
笔记：\path{reading_notes/hall_fluid_review.md}。\par\smallskip
\paragraph{Kaptanoglu2020（2020）}\label{paper:Kaptanoglu2020}
Two-temperature effects in Hall-MHD simulations of the HIT-SI experiment。11月12日新增全文核读；本次修订不重复计数；主版本35个PDF物理页。
OSTI author accepted manuscript, 35 pages: article pp1–18 and enlarged/repeated figures pp19–35. Unreviewed November 2020 erratum exists.。\href{run:../../papers/Kaptanoglu2020_accepted.pdf}{打开本地原文}。
笔记：\path{reading_notes/hall_fluid_review.md}。\par\smallskip
\paragraph{TenBarge2019MMS（2019）}\label{paper:TenBarge2019MMS}
An Extended MHD Study of the 16 October 2015 MMS Diffusion Region Crossing。11月12日新增全文核读；本次修订不重复计数；主版本17个PDF物理页。
arXiv:1903.06605v1；正式HTML方法/讨论补查，正式PDF403。\href{run:../../papers/TenBarge2019MMS.pdf}{打开本地原文}。
笔记：\path{reading_notes/closure_review.md}。\par\smallskip
\paragraph{F674（2015）}\label{paper:F674}
Study of energy conversion and partitioning in the magnetic reconnection layer of a laboratory plasma。11月12日新增全文核读；本次修订不重复计数；主版本11个PDF物理页。
publisher\_version。\href{run:../../papers/F674_published.pdf}{打开本地原文}。
笔记：\path{reading_notes/experiment_review.md}。\par\smallskip
\paragraph{Yoo2017Heating（2017）}\label{paper:Yoo2017Heating}
Electron heating and energy inventory during asymmetric reconnection in a laboratory plasma。11月12日新增全文核读；本次修订不重复计数；主版本18个PDF物理页。
publisher\_version。\href{run:../../papers/Yoo2017Heating.pdf}{打开本地原文}。
笔记：\path{reading_notes/experiment_review.md}。\par\smallskip
\paragraph{Yoo2013（2013）}\label{paper:Yoo2013}
Observation of Ion Acceleration and Heating during Collisionless Magnetic Reconnection in a Laboratory Plasma。11月12日完整重读；本次修订不重复计数；主版本5个PDF物理页。
publisher\_version。\href{run:../../papers/Yoo2013.pdf}{打开本地原文}。
笔记：\path{reading_notes/experiment_review.md}。\par\smallskip
\paragraph{L272（2021）}\label{paper:L272}
How Hall electric fields intrinsically chaotize and heat ions during collisionless magnetic reconnection。11月12日完整重读；本次修订不重复计数；主版本10个PDF物理页。
publisher\_version。\href{run:../../papers/L272_published.pdf}{打开本地原文}。
笔记：\path{reading_notes/experiment_review.md}。\par\smallskip
\paragraph{E512（2011）}\label{paper:E512}
Reduced fluid-kinetic equations for low-frequency dynamics, magnetic reconnection, and electron heating in low-beta plasmas。11月12日完整重读；本次修订不重复计数；主版本28个PDF物理页。
arXiv:1104.4622v2。\href{run:../../papers/E512.pdf}{打开本地原文}。
笔记：\path{reading_notes/closure_review.md}。\par\smallskip
\paragraph{Yoon2023PhaseSpace（2023）}\label{paper:Yoon2023PhaseSpace}
Phase-space Analysis of Ordered and Disordered Nonthermal Ion Energization during Magnetic Reconnection。11月12日新增全文核读；本次修订不重复计数；主版本9个PDF物理页。
publisher\_version。\href{run:../../papers/Yoon2023PhaseSpace.pdf}{打开本地原文}。
笔记：\path{reading_notes/experiment_review.md}。\par\smallskip
\paragraph{F679（2014）}\label{paper:F679}
Bulk ion acceleration and particle heating during magnetic reconnection in a laboratory plasma。11月12日新增全文核读；本次修订不重复计数；主版本13个PDF物理页。
publisher\_version。\href{run:../../papers/F679.pdf}{打开本地原文}。
笔记：\path{reading_notes/experiment_review.md}。\par\smallskip
\paragraph{L263（2019）}\label{paper:L263}
Fast Recursive Reconnection and the Hall Effect: Hall-MHD Simulations。复用原库已核读档案；未声称本次重新全文阅读；主版本13个PDF物理页。
原库登记版本。\href{run:../../papers/L263.pdf}{打开本地原文}。
笔记：\path{research/reading_notes/verified.json}。\par\smallskip
\paragraph{Stanier2013（2013）}\label{paper:Stanier2013}
Two-fluid simulations of driven reconnection in the mega-ampere spherical tokamak。复用原库已核读档案；未声称本次重新全文阅读；主版本12个PDF物理页。
原库登记版本。\href{run:../../papers/Stanier2013.pdf}{打开本地原文}。
笔记：\path{research/reading_notes/verified.json}。\par\smallskip
\paragraph{L247（2019）}\label{paper:L247}
High-order two-fluid plasma solver for direct numerical simulations of plasma flows with full transport phenomena。11月12日新增全文核读；本次修订不重复计数；主版本44个PDF物理页。
arXiv1812.11237v1 2018-12-28, accepted for 2019 article; catalog previously used preprint year2018。\href{run:../../papers/L247.pdf}{打开本地原文}。
笔记：\path{reading_notes/hall_fluid_review.md}。\par\smallskip
\paragraph{L244（2018）}\label{paper:L244}
A reduced Landau-gyrofluid model for magnetic reconnection driven by electron inertia。11月12日完整重读；本次修订不重复计数；主版本27个PDF物理页。
HAL hal-02049658v1，含封面。\href{run:../../papers/L244.pdf}{打开本地原文}。
笔记：\path{reading_notes/closure_review.md}。\par\smallskip
\paragraph{Yoon2018Heating（2018）}\label{paper:Yoon2018Heating}
Fast Ion Heating in Transient Collisionless Magnetic Reconnection via an Intrinsic Stochastic Mechanism。11月12日新增全文核读；本次修订不重复计数；主版本7个PDF物理页。
publisher\_version。\href{run:../../papers/Yoon2018Heating.pdf}{打开本地原文}。
笔记：\path{reading_notes/experiment_review.md}。\par\smallskip
\paragraph{Finelli2021（2021）}\label{paper:Finelli2021}
Bridging hybrid- and full-kinetic models with Landau-fluid electrons. I. 2D magnetic reconnection。11月12日新增全文核读；本次修订不重复计数；主版本17个PDF物理页。
arXiv:2101.00040v2，17页8图；正式PDF403。\href{run:../../papers/Finelli2021.pdf}{打开本地原文}。
笔记：\path{reading_notes/closure_review.md}。\par\smallskip
\paragraph{Numata2015Heating（2015）}\label{paper:Numata2015Heating}
Ion and electron heating during magnetic reconnection in weakly collisional plasmas。11月12日新增全文核读；本次修订不重复计数；主版本18个PDF物理页。
arXiv:1406.6456v2；2014在线发表、2015卷期。\href{run:../../papers/Numata2015Heating.pdf}{打开本地原文}。
笔记：\path{reading_notes/closure_review.md}。\par\smallskip
\paragraph{Ng2020（2020）}\label{paper:Ng2020}
An improved ten-moment closure for reconnection and instabilities。复用原库已核读档案；未声称本次重新全文阅读；主版本8个PDF物理页。
原库登记版本。\href{run:../../papers/Ng2020.pdf}{打开本地原文}。
笔记：\path{research/reading_notes/verified.json}。\par\smallskip
\paragraph{Closure2026（2026）}\label{paper:Closure2026}
Capturing Secondary Kinetic Instabilities in Three‐Dimensional Dayside Reconnection Using an Improved Gradient‐Based Closure。复用原库已核读档案；未声称本次重新全文阅读；主版本23个PDF物理页。
原库登记版本。\href{run:../../papers/Closure2026.pdf}{打开本地原文}。
笔记：\path{research/reading_notes/verified.json}。\par\smallskip
\paragraph{F635（2026）}\label{paper:F635}
Counter-streaming heat-flux closure for electron-only collisionless magnetic reconnection。11月12日完整重读；本次修订不重复计数；主版本10个PDF物理页。
原库登记版本。\href{run:../../papers/F635.pdf}{打开本地原文}。
笔记：\path{reading_notes/F635_note.json}。\par\smallskip
\paragraph{L164（2010）}\label{paper:L164}
Slow shocks and conduction fronts from Petschek reconnection of skewed magnetic fields: two-fluid effects。11月12日新增全文核读；本次修订不重复计数；主版本45个PDF物理页。
arXiv:1006.0441v1 (2010-06-02); local regenerated title page says draft May 29 2018; not a new publication。\href{run:../../papers/L164.pdf}{打开本地原文}。
笔记：\path{reading_notes/hall_fluid_review.md}。\par\smallskip
\paragraph{F743（2019）}\label{paper:F743}
Kinetic Verification of the Stochastic Ion Heating Mechanism in Collisionless Magnetic Reconnection。11月12日新增全文核读；本次修订不重复计数；主版本5个PDF物理页。
publisher\_version。\href{run:../../papers/F743.pdf}{打开本地原文}。
笔记：\path{reading_notes/experiment_review.md}。\par\smallskip
\paragraph{Loureiro2013Heating（2013）}\label{paper:Loureiro2013Heating}
Fast collisionless reconnection and electron heating in strongly magnetized plasmas。11月12日新增全文核读；本次修订不重复计数；主版本6个PDF物理页。
arXiv:1301.0338v1；排版另印Dated May 5,2018，采用存档v1标识。\href{run:../../papers/Loureiro2013Heating.pdf}{打开本地原文}。
笔记：\path{reading_notes/closure_review.md}。\par\smallskip
\paragraph{Che2014Buneman（2014）}\label{paper:Che2014Buneman}
Two-Fluid Description of Wave-Particle Interactions in Strong Buneman Turbulence。11月12日新增全文核读；本次修订不重复计数；主版本26个PDF物理页。
arXiv:1403.4552v2。\href{run:../../papers/Che2014Buneman.pdf}{打开本地原文}。
笔记：\path{reading_notes/closure_review.md}。\par\smallskip
\paragraph{Drake2009Pickup（2009）}\label{paper:Drake2009Pickup}
Ion heating resulting from pickup in magnetic reconnection exhausts。11月13日新增全文核读；主版本14个PDF物理页。
JGR正式版，机构作者网页副本，14页含附录A。\href{run:../../papers/Drake2009Pickup.pdf}{打开本地原文}。
笔记：\path{reading_notes/Drake2009Pickup_note.json}。\par\smallskip
\paragraph{Tanabe2015MAST（2015）}\label{paper:Tanabe2015MAST}
Electron and Ion Heating Characteristics during Magnetic Reconnection in the MAST Spherical Tokamak。11月12日新增全文核读；本次修订不重复计数；主版本5个PDF物理页。
publisher\_version。\href{run:../../papers/Tanabe2015MAST.pdf}{打开本地原文}。
笔记：\path{reading_notes/experiment_review.md}。\par\smallskip
\paragraph{TS6FEC2025（2025）}\label{paper:TS6FEC2025}
The scaling of the ion heating and electrostatic potential in spherical tokamak。11月12日新增全文核读；本次修订不重复计数；主版本8个PDF物理页。
Watermarked PREPRINT—IAEA-FEC2025 manuscript; no peer-reviewed journal version verified。\href{run:../../papers/TS6FEC2025.pdf}{打开本地原文}。
笔记：\path{reading_notes/hall_fluid_review.md}。\par\smallskip

\chapter{选定代数核验与阅读边界}
11月12日用SymPy核查四项教学恒等式：Hall总功和分物种功相消、摩擦流能损失与$\eta J^2$一致、十矩局地松弛的零迹、等密度两束的中心二阶矩。脚本\path{scripts/check_ion_heating_algebra.py}和结果\path{research/ion_heating/algebra_checks.json}可复查。11月13日另补带符号电场梯度轨道方程及拾取温度迹，共七个检查项（Hall功含总功与物种和两项）。它们只验证选定代数，不证明所有公式、闭合或模拟。

11月12日批次由三条阅读分支完成，主报告做跨分支定义及证据一致性核对；11月13日修订对旧文作定向正文/图页复查，并新增Drake2009Pickup全文14页，详见revision\_manifest.json与reading\_notes/Drake2009Pickup\_note.json。已核读不等于作者代码、独立补充和原始数据都已运行核验。阅读笔记明确保留版本差异、未报告参数和未取得资料。原仓库122篇基线不是本轮重新逐篇阅读了122篇，更不是通读全部原文库；本轮新增与重读分别计数。

主报告源文件为本目录\path{report.tex}，证据索引/书目由\path{scripts/build_ion_heating_report.py}生成。全局源章节同步加入研究线和问题总结；三卷由既有构建脚本生成。本专题的检索原始命中单独归档，未将大量未筛选元数据混入全局科学证据目录。

\backmatter
\bibliographystyle{unsrtnat}
\bibliography{references}
\end{document}
